\section{Complete reciprocal charts and scaled logarithms}

Put $K=\Q(\rho)$, $\rho^2=2$, $d=3-\rho$, $D=d^2$, and
$E:y^2=x^3-Dx$. The four root points retain the charts
\[
 x=h^2,\qquad j^2=x^2-D,\qquad y=hj.
\]
Their $h,j,x$ are nonzero. Indeed $h=0$ would give $j^2=-D$, impossible
under a real embedding. If $j=0$, then $x=\pm d$, whose norm is $7$,
whereas the norm of a square is a rational square. The inherited group
description identifies these charts with finite points of $2E(K)=\Z P$.

\subsection{Primary charts}
For $H,U_1,\ldots,U_4$ use instead
\begin{equation}\label{eq:reciprocal-chart}
 f=g^2,\qquad l^2=1-Df^2.
\end{equation}
For $g\ne0$ this represents $x=1/g^2$, $y=l/g^3$, with old coordinates
$h=1/g$, $j=l/g^2$. Conversely $g=1/h$, $l=j/h^2$. The orientations agree.
When $g=0$, $l=\pm1$; both pairs represent the identity $O$.

The coordinate $l$ never vanishes: otherwise $dg^2=\pm1$, so
$7\Norm(g)^2=1$, impossible in $\Q$. Thus $z=-g/l$ is defined everywhere
on the chart, agreeing with $-x/y$ on its finite locus and taking value zero
at $O$. Similarly $f$ extends $1/x$ with value zero at $O$.

Define the shared inverse parameter and primary ratios by
\[
 g_Hk_H=l_H,\qquad l_iL_i=k_Hg_i.
\]
The first equation forces $g_H\ne0$, so $H$ remains finite and nonidentity.
The second gives $L_i=z(U_i)/z(H)$, including $L_i=0$ when $U_i=O$.
Five charts save ten rational coordinates; the shared $k_H$ costs two.

\subsection{A reciprocal difference formula}
Let $R=(x,hj)$ be a retained finite root, with $x=h^2$ and $j^2=x^2-D$.
For $U=(g,l,f)$ as above introduce
\begin{equation}\label{eq:complete-difference}
 \begin{aligned}
 F&=x+Df, &F\tau&=hl-gj,\\
 f_V&=\tau^2, &\tau M&=-g(x+Df_V).
 \end{aligned}
\end{equation}
Conceptually put $k_V=(M-j)/h$ and
$V=(1/f_V,k_V/f_V)$. All displayed equations are quadratic in the
materialized coordinates.

The first denominator is totally positive: $x=h^2>0$ at both real
embeddings, $D>0$, and $f=g^2\ge0$. The equations also force $\tau\ne0$.
If $\tau=0$, then $hl=gj$, and squaring gives
$(1-xf)(x+Df)=0$. Thus $xf=1$ and $g\ne0$, whereas the last equation
would say $0=-gx$. Hence $f_V\ne0$, and $M$ is uniquely determined.

For the generic identity set
\[
 A=hl-gj,\qquad B=hl+gj,\qquad \delta=1-xg^2.
\]
The exact relation
\[
 AB-\delta F
 =x(l^2-1+Dg^4)-g^2(j^2-x^2+D)=0
\]
shows, when $g\delta\ne0$, that the usual subtraction abscissa
$B^2/\delta^2$ equals $F^2/A^2=1/\tau^2$.
Substitution in the usual ordinate formula gives
\[
 \frac{y(U-R)}{x(U-R)}
 =-\frac jh-\frac{g(x+D\tau^2)}{h\tau}
 =\frac{M-j}{h}.
\]
The independent checker verifies the cleared abscissa and ordinate
identities modulo the two chart equations, as well as the resulting curve
equation. Thus the orientation is fixed, not selected up to sign.

The exceptional cases are explicit. If $U=O$, then
\[
 \tau=l/h,\quad f_V=1/x,\quad M=0,\quad k_V=-j/h,
\]
so $V=-R$. If $U=-R$, write $\epsilon=gh\in\{1,-1\}$ and
$A_R=x^2+D$, $y_R=hj$. Equations~\eqref{eq:complete-difference} give
\[
 \tau=-2\epsilon y_R/A_R,\quad
 x(V)=A_R^2/(4y_R^2),\quad
 k_V=-\frac{x^4-6Dx^2+D^2}{2y_RA_R}.
\]
These are exactly the coordinates of $-2R$. This covers tangency, since
$y_R\ne0$ and $A_R>0$. If $U=R$, the first numerator is zero, so
$\tau=0$, already proved impossible. The identity output remains excluded.
These cases exhaust all pairs. Consequently every output is exactly the
finite nonidentity point $U-R$ in $2E(K)$, with nonzero $k_V$.

\subsection{The scalar difference with a fixed target}
The scalar root $S$ retains its affine $h_S,j_S,x_S$ chart. With fixed
$P=(x_P,h_Pj_P)$ use
\[
 \begin{aligned}
 (x_Px_S+D)\tau&=h_Pj_S-h_Sj_P,\quad f_R=\tau^2,\\
 w&=h_S\tau,\quad wM=-(x_P+Df_R).
 \end{aligned}
\]
Then $k_R=(M-j_P)/h_P$ is affine and the output is $S-P$.
This follows by substituting $g=1/h_S$, $l=j_S/x_S$ in the preceding
formula. It includes $S=-P$ and excludes $S=P$. The anchor still forces
the same large nonzero scalar index $a-1$. All six difference guards
therefore disappear, and the old $2$-adic difference-component lemma
is no longer required.

\subsection{Scale logarithms by the root coordinate}
The alternate logarithm remains
$\ell(V)=-G_9(Df_V^2)/k_V$. Instead of materializing $k_V$ at a variable
target, store $\eta=\ell(V)/h$ by
\[
 (M-j)\eta=-G_9(Df_V^2).
\]
The final quadratic stage of $G_9$ can be substituted directly here.
The actual ratio consumer is
\[
 L_i'=\overline{\frac{h_i\eta_i}{h_P\eta_0}}.
\]
Its numerator is quadratic; its denominator is affine because $h_P$ is
fixed. Thus the four $k_{V_i}$ pairs are unnecessary, saving eight rational
coordinates. Conjugation still applies to the entire ratio.

All denominators are nonzero. In particular the inherited nonvanishing of
$G_9$ can also be seen from the source's elementary real argument:
if $G_J(w)=1+\sum_{j=1}^J c_jw^j$ with the retained coefficients, then
$c_j<0$ and $G_J(1)=-(4J-1)c_J/2>0$. A finite point on these charts has
$0<D/x^2<1$, so $G_J(D/x^2)>0$ at a real embedding.

\subsection{Soundness of the new identity-input case}
Every old generic witness tuple extends to the new point equations.
The converse reconstructs the same mathematical functions on a larger
domain: tangency is allowed, and each $U_i$ may equal $O$. This is not
a bijection with the old guarded witness set.

The actual local-depth consumer is the subsection ``Local soundness'' in\\
\nolinkurl{docs/quartic-523-540/addition.tex}. Write
\[
 H=[h]P,\quad U_i=[u_i]P,\quad R_i=[n_i]P,\qquad
 r_i=u_i/h,\quad r_i'=(u_i-n_i)/(h-1).
\]
Here $h\ne0,1$, every $n_i\ne0$, and $u_i\ne n_i$; none requires
$u_i\ne0$. If $U_i=O$, then $L_i=r_i=0$. The order-40 cross test gives
$v(L_i')\ge40$. Both alternate points remain nonidentity formal points,
so their unchanged relative logarithm estimate gives $v(r_i')\ge40$.
Therefore
\[
 n_i=-(h-1)r_i'\equiv0=L_i\pmod{p^{40}}.
\]
This is exactly the conclusion used by the beta and Lucas tests. No
undefined ratio $\log_E(O)/z(O)$ is used. For finite $U_i$ the old proof
applies unchanged; at inert primes conjugation preserves the valuation.
The subsequent height arguments use only the unchanged root labels and
these congruences. Completeness can retain its old finite generic choices.

The total geometric saving is $10-2+6+8=22$, giving 428 variables in
isolation. The accompanying full proof text records every domain and
the precise source application; exact checks cover each exceptional case.
